\documentclass[exercice,a4]{thedocument}%
\usepackage{texmaths-tikz}%
\usepackage{pstricks}%
\usepackage{pst-bezier}%
\usepackage{pstricks-add}%
\usepackage{pst-eucl}%
\usepackage{pst-func}%
\usepackage{pst-tree}%
\usepackage{pst-circ}%
\usepackage{tikz}%
\usepackage{pgfplots}%
\usepackage{tkz-tab}%
\usepackage{tikzpagenodes}%
\usepackage{tabularray}%
\usepackage{tabularx}%
\usepackage{multirow}%
\usepackage{multicol}%
\usepackage{scratch3}%
\usepackage{enumitem}%
\usepackage[np]{numprint}%
\usepackage{pas-tableur}%
\usepackage{listings}%
\usepackage{graphicx}%
\usepackage{adjustbox}%
\usepackage{needspace}%
\usepackage{xcolor}%
\usepackage{eurosym}%
\usepackage{diagbox}%
\usepackage{colortbl}%
\usepackage{multido}%
\usepackage{icomma}%
\usepackage[french]{babel}%
\startdatakeys
    \source{26GENMATAA1}
    \BO{2026}
    \cycle{4}
    \competencesDuSocle{CA,RA,CO}
    \objectifsApprentissage{C4A1aa,C4A2ad,C4B2cb,C4D2aa,C4D3ac}
    \motsCles{Pythagore,droites parallèles,Thalès,cône,volume}
\enddatakeys
\definecolor{gris}{gray}{0.5}%
\providecommand{\wideparen}[1]{\overset{\frown}{#1}}%
\begin{document}%
\providecommand{\degres}{^\circ}%
\emph{La figure ci-dessous n'est pas à l'échelle.}

\begin{minipage}{0.48\linewidth}
	\begin{itemize}
		\item O, A et B sont alignés.
		\item O, D et C sont alignés.
		\item OD $= 8,2$ cm
		\item AD $= 1,8$ cm
		\item BC $= 4,5$ cm
	\end{itemize}
\end{minipage}\hfill
\begin{minipage}{0.48\linewidth}
	% Début PSTricks
	\psset{unit=1cm}
	\begin{pspicture}(-1,-0.5)(3.6,5)
		%\psgrid
		\pspolygon(0,0)(0,5)(2.6,5)%ABCO
		\psline(0,3)(1.6,3)%AD
		\psframe(0,3)(0.3,2.7,3)
		\psframe(0,5)(0.3,4.7)
		\uput[d](0,0){O} \uput[l](0,3){A} \uput[l](0,5){B} \uput[r](2.6,5){C} \uput[r](1.6,3){D}
	\end{pspicture}
	%Fin  PSTricks
\end{minipage}
%\begin{tikzpicture}[scale=1.05]
%% Points based on description: O at bottom, A,B aligned ; D,C aligned
%\coordinate (O) at (0,0);
%\coordinate (A) at (0.7,3.6);
%\coordinate (D) at (1.55,3.6);
%\coordinate (B) at (1.0,5.0);
%\coordinate (C) at (2.5,5.0);
%
%\draw (O) -- (B);
%\draw (O) -- (C);
%\draw (A) -- (D);
%\draw (B) -- (C);
%
%% right angle marks
%\draw ($(A)+(0.18,0)$) -- ($(A)+(0.18,0.18)$) -- ($(A)+(0,0.18)$);
%\draw ($(B)+(0.18,0)$) -- ($(B)+(0.18,0.18)$) -- ($(B)+(0,0.18)$);
%
%\node[left] at (O) {$O$};
%\node[left] at (A) {$A$};
%\node[right] at (D) {$D$};
%\node[above left] at (B) {$B$};
%\node[above right] at (C) {$C$};
%\end{tikzpicture}

\medskip

\begin{enumerate}
	\item Montrer que la longueur du segment [OA] est égale à 8 cm.
	\item Justifier que les droites (BC) et (AD) sont parallèles.
	\item Calculer la longueur du segment [OB].
	\item Une entreprise souhaite fabriquer des gobelets. Un gobelet (grisé sur le schéma ci-dessous) a la forme d'un tronc de cône (cône coupé par un plan parallèle à sa base).

	\medskip

	\begin{minipage}{0.48\linewidth}
		On reprend les données précédentes :

		\begin{itemize}
			\item O, A, B sont alignés
			\item O, D, C sont alignés
			\item OD $= 8,2$ cm
			\item AD $= 1,8$ cm
			\item BC $= 4,5$ cm
		\end{itemize}
	\end{minipage}\hfill
	\begin{minipage}{0.48\linewidth}
		\emph{La figure ci-dessous n'est pas à l'échelle.}

		\medskip

		%Début Tikz
%		\begin{center}
%		\begin{tikzpicture}[scale=0.9]
%		% big cone with frustum shaded
%		\coordinate (O) at (0,0);
%		\coordinate (Atop) at (-1.0,3.2); % small ellipse left
%		\coordinate (Btop) at (-2.4,4.6); % large ellipse left
%
%		% Outer cone lines
%		\draw[dashed] (O) -- (-2.4,4.6);
%		\draw[dashed] (O) -- (2.4,4.6);
%		\draw[dashed] (O) -- (-1.0,3.2);
%		\draw[dashed] (O) -- (1.0,3.2);
%
%		% Large ellipse (top, cup rim)
%		\draw (-2.4,4.6) ellipse (2.4 and 0.5);
%		% small ellipse (base of cup)
%		\draw (-1.0,3.2) ellipse (1.0 and 0.22);
%
%		% Shaded frustum sides
%		\filldraw[fill=gray!25, draw=black] (-1.0,3.2) -- (-2.4,4.1) arc (180:0:2.4 and 0.5) -- (1.0,3.2) arc(0:180:1.0 and 0.22) -- cycle;
%
%		\node[right] at (2.4,4.6) {$C$};
%		\node[left] at (-2.4,4.6) {$B$};
%		\node[right] at (1.0,3.2) {$D$};
%		\node[left] at (-1.0,3.2) {$A$};
%		\node[below] at (O) {$O$};
%
%		\draw[->, thick] (3.6,3.8) -- (1.6,3.9);
%		\node[right] at (3.6,3.8) {Le gobelet};
%		\end{tikzpicture}
%		\end{center}
		%Fin Tikz
		%Début PSTricks
		\begin{pspicture}(-2.6,-0.3)(2.6,5)
			%\psgrid
			\psellipse(0,4.8)(2.3,0.3)
			\psline[linestyle=dashed](-2.34,4.8)(2.34,4.8)
			\psellipticarc(0,2.3)(1.15,0.2){180}{360}
			\psellipticarc[linestyle=dashed](0,2.3)(1.15,0.2){0}{180}
			\psline[linestyle=dashed](0,0)(0,4.8)%OB
			%\psline[linestyle=dashed](-2.3,5)(2.3,5)
			\psline[linestyle=dashed](-1.15,2.3)(1.15,2.3)
			\psline(-2.34,4.8)(-1.13,2.3)
			\psline(2.34,4.8)(1.13,2.3)
			\psframe(0,4.8)(0.25,4.55)\psframe(0,2.3)(0.25,2.55)
			\psline[linestyle=dashed](-1.15,2.3)(0,0)(1.15,2.3)
			\uput[d](0,0){O} \uput[ul](0,2.3){A} \uput[ul](0,4.8){B} \uput[r](2.3,4.8){C} \uput[r](1.15,2.3){D}
			\rput(0,3.6){Le gobelet}
		\end{pspicture}
		%Fin PSTricks
	\end{minipage}


	\begin{center}
		\fbox{\parbox{0.7\textwidth}{\textbf{Rappel :}\\
				Volume d'un cône de révolution :\quad $V = \pi \times R^2 \times H \div 3$,\\ où $R$ est le rayon de la base et $H$ est la hauteur du cône.}}
	\end{center}

	\begin{enumerate}
		\item Calculer le volume du grand cône de hauteur [OB] en cm$^3$, arrondi à l'unité.
		\item Calculer le volume du gobelet, en cm$^3$, arrondi à l'unité.
	\end{enumerate}
\end{enumerate}

\newpage
\medskip
\showthesource
\end{document}
